\section{An arithmetic consumer with all cyclic layers and parity obstructions}
\label{sec:layeredcoprime}
For fixed \(k\geq2\), form the directed graph on
\((\mathbb Z/k\mathbb Z)\times(\mathbb Z/n\mathbb Z)\)
with arcs
\[
 (i,x)\longrightarrow(i+1,y)
             \quad\Longleftrightarrow\quad \gcd(y-x,n)=1 .
\]
Let \(H_k(n)\) count directed Hamilton cycles modulo cyclic
rotation, retaining their orientation. This convention is
important: for \(k=2\) opposite orientations are both counted.
The graph is \(\varphi(n)\)-regular on \(kn\) vertices.

For odd \(s\), including \(s=1\), put
\(\lambda_d=\mu(d)/\varphi(d)\) and define
\[
 {\cal D}_k(s)=
 \prod_{\substack{d\mid s\\d>1}}
 \left(\frac{1-\lambda_d^k}
              {(1-\lambda_d^2)^{k/2}}\right)^{\varphi(d)},
 \qquad
 {\cal E}_k(s)=
 \prod_{\substack{d\mid s\\d>1}}
 \left(\frac{1-\lambda_d^{2k}}
              {(1-\lambda_d^2)^k}\right)^{\varphi(d)} .
\]
Empty products equal one; factors with \(\mu(d)=0\) equal one.
All displayed factors are positive.

\begin{theorem}[Complete relative counts for layered coprime circles]
\label{thm:layeredcoprime}
Fix \(k\geq2\). Uniformly as \(n\to\infty\):
\begin{enumerate}
\item If \(n\) is odd, then
\[
 H_k(n)=
 \frac{(n!)^k}{n\,n^{kn}}\varphi(n)^{kn}{\cal D}_k(n)
                  \left(1+O_k(n^{-1/10})\right).
\]
Equivalently its leading prefactor is
\((2\pi)^{k/2}n^{k/2-1}{\cal D}_k(n)\)
multiplying \((\varphi(n)/e)^{kn}\).
\item If both \(n\) and \(k\) are even, then \(H_k(n)=0\) exactly.
\item If \(n=2m\) is even and \(k\) is odd, write
\(n=2^a s\) with \(s\) odd. Then
\[
 H_k(n)=
 \frac{(m!)^{2k}}{m\,m^{2km}}\varphi(n)^{kn}{\cal E}_k(s)
                   \left(1+O_k(n^{-1/10})\right).
\]
Equivalently its leading prefactor is
\((2\pi)^k(n/2)^{k-1}{\cal E}_k(s)\)
multiplying \((\varphi(n)/e)^{kn}\).
\end{enumerate}
There is no restriction on the number of distinct prime
divisors, on their sizes, or on the two-adic valuation.
The constants may depend on the fixed \(k\).
\end{theorem}
\begin{proof}
For odd \(n\), every normalized cross block is the same
symmetric doubly stochastic unitary kernel \(T\).
Its nontrivial eigenvalues are \(\lambda_d\), each with
multiplicity \(\varphi(d)\), for \(d\mid n\), \(d>1\).
They have absolute value at most \(1/2\).
The entry parameter is \(A=O(\log\log(3n))\).
Theorem~\ref{thm:periodiccomparison} applies with \(m=n\);
its full determinant quotient is precisely \({\cal D}_k(n)\).
Restoring the common edge-weight factor
\(\varphi(n)^{kn}\) proves the first formula.

When \(n\) and \(k\) are both even, every arc preserves
\(x+i\pmod2\). Both invariant classes are nonempty.
No directed Hamilton cycle can cover both, proving exact
vanishing, rather than a failure of an asymptotic estimate.

Now suppose \(n=2m\) and \(k\) is odd.
Every arc flips the parity of the coordinate \(x\).
The map on layer/parity states
\((i,\epsilon)\mapsto(i+1,1-\epsilon)\) is a single cycle
of length \(2k\). Thus the actual full graph has \(2k\)
cyclic layers, each of size \(m\).
Their cross blocks alternate between \(U\) and \(U^\top\),
where \(U\) is the normalized even-to-odd unitary block.

By the Chinese remainder identification modulo \(2^a\)
and \(s\), row and column permutations identify \(U\) with
the tensor product of the unitary kernel on
\(\mathbb Z/s\mathbb Z\) and the constant stochastic
matrix of order \(2^{a-1}\).
For \(s=1\) the former kernel is the \(1\)-by-\(1\) matrix \((1)\).
Consequently the squared nontrivial singular values of \(U\)
are \(\lambda_d^2\) with multiplicity \(\varphi(d)\),
for \(d\mid s\), \(d>1\), together with zero values.
Their maximum is at most \(1/4\), uniformly in \(a,s\).

Write \(B=U-J_m/m\). The cyclic product of the \(2k\)
centered cross blocks is \((BB^\top)^k\).
The full numerator is \(\det(I-(BB^\top)^k)\);
the product of square-root denominators is
\(\det(I-BB^\top)^k\).
Their quotient is \({\cal E}_k(s)\).
The density parameter satisfies
\[
 m\max U_{ij}=\frac{m}{\varphi(n)}
            =\frac{s}{\varphi(s)}
               =O(\log\log(3n)).
\]
Apply Theorem~\ref{thm:periodiccomparison} with \(2k\) layers
and layer size \(m\), and restore \(\varphi(n)^{kn}\).
Its error is at most
\(O_k(n^{-1/9}\exp(C_k\log\log(3n)))\), hence
\(O_k(n^{-1/10})\).
The same bound applies in the odd case.
Finally Stirling's formula proves the two equivalent
prefactor expressions. All parity modes, orientations,
and layer-size factorials have been included.
\end{proof}

\begin{proposition}[The uncancelled arithmetic interaction for higher layers]
\label{prop:layeredconstants}
For odd \(s\), define, for integers \(b\geq2\),
\[
 F_b(s)=\prod_{p\mid s}
              \left(1+\frac{(-1)^b}{(p-1)^{b-1}}\right)-1 .
\]
Then \(F_2(s)=s/\varphi(s)-1\), and exactly
\[
 \log{\cal D}_k(s)
 =\sum_{j\geq1}
       \left(\frac{k}{2j}F_{2j}(s)-\frac1jF_{kj}(s)\right),
 \qquad
 \log{\cal E}_k(s)
 =\sum_{j\geq1}
       \left(\frac{k}{j}F_{2j}(s)-\frac1jF_{2kj}(s)\right).
\]
In particular \({\cal D}_2(s)=1\) exactly.
For every fixed \(k\geq3\),
\[
 \log{\cal D}_k(s)
      =\frac k2\left(\frac{s}{\varphi(s)}-1\right)+O_k(1),
 \qquad
 \log{\cal E}_k(s)
      =k\left(\frac{s}{\varphi(s)}-1\right)+O_k(1),
\]
uniformly over all odd \(s\).
After subtracting those respective leading terms, the
remaining series converge absolutely and uniformly.
The corresponding normalized constants have finite positive
limits along odd primorials.
\end{proposition}
\begin{proof}
For \(b\geq2\), the exact multiplicative identity is
\[
 \sum_{\substack{d\mid s\\d>1}}\varphi(d)\lambda_d^b
       =\sum_{\substack{d\mid s\\d>1}}
                   \frac{\mu(d)^b}{\varphi(d)^{b-1}}
       =F_b(s).
\]
Expand each logarithm in the definitions of
\({\cal D}_k,{\cal E}_k\), retaining the numerator and
denominator series with their signs. Since every
nontrivial \(|\lambda_d|\leq1/2\), these expansions
converge for each \(s\), and give the displayed formulas.
For \(k=2\) they cancel term by term.

Uniformity after the leading term is also explicit.
For \(b\geq3\),
\[
 |F_b(s)|
 \leq\exp\left(\sum_{p\ {\rm odd}}(p-1)^{-(b-1)}\right)-1
 \leq\exp\left(2^{-(b-1)}\zeta(b-1)\right)-1
 \leq C\,2^{-b}.
\]
The constant \(C\) is absolute, as \(\zeta(b-1)\leq2\).
This gives a summable bound for every remaining term.
In \({\cal D}_k\), \(k\geq3\), the only \(b=2\) term
is \(kF_2/2\); in \({\cal E}_k\) it is \(kF_2\).
All other terms have \(b\geq3\).
Each infinite-prime product with \(b\geq3\) converges
absolutely along primorials. Dominated convergence now
proves the asserted normalized limits; exponentiation
makes them positive and finite.
\end{proof}

\begin{remark}[The mechanism and its limits]
For a single symmetric pair of cross blocks the quadratic
spectral interaction cancels exactly.
For three or more layers it survives as the explicit
\(s/\varphi(s)-1\) contribution, with a uniformly controlled
higher-order prime interaction.
Thus the method can distinguish exact cancellation from
a nonzero arithmetic constant, rather than just relocating
an unestimated signed term.
It proves relative counts on the entire modulus domain in
this theorem. It makes no assertion about integer prime
gaps, RH, or a zero-free exponent.
The fixed-layer restriction and the stated diffuseness
range remain essential hypotheses of the comparison used.
\end{remark}
